%==============================================================================
\section{State of the art}\label{sec:sota}
%==============================================================================

This section reviews the SSP literature. It is organised around one question:
how strong a lower bound can be obtained, and by what means. That question
organises the exact literature completely, and it also explains why the heuristic
literature looks the way it does.

\subsection{Origins: the loading rule, and the hardness of ordering}\label{ssec:sota-origins}

\citet{Tang1988} posed the problem and settled its inner level. For a fixed job
order the optimal loading is the Keep Tool Needed Soonest rule, proved optimal and
computable in $O(n\lvert T\rvert)$ time (Proposition~\ref{prop:ktns}). Two further
elements of that paper are used directly in this report. The first is the pairwise
bound: if jobs $i$ and $j$ run consecutively, at least
$\max(0,\lvert T_i\cup T_j\rvert-b)$ tools must be switched between them, because
the magazine holds $b$ tools but must serve $T_i$ and then $T_j$. The second is
their observation that the shortest Hamiltonian path under those pairwise weights
is a lower bound on the optimum. Both reappear in the Benders master of
Section~\ref{ssec:bbc}.

\citet{Crama1994} made the difficulty precise. They proved the SSP NP-hard for every
fixed capacity $b\ge2$, by showing that finding a minimum-length travelling-salesman
path in the edge graph of an arbitrary graph can be written as an SSP instance with
$b=2$, and then extending the construction to larger $b$ by padding the incidence
matrix. The same paper fixes the seven modelling assumptions that have been standard
in the field since: the magazine is always full, each tool occupies one slot, switch
times are equal for all tools, tools are changed one at a time, tool requirements are
fixed in advance, tools do not wear out, and the job list is known.

That paper also gives a second, linear-programming proof of KTNS optimality through
interval matrices, and its treatment of the non-uniform case is relevant to
Section~\ref{ssec:bbc}. When tools have different setup times, the greedy argument
breaks: \citet{Crama1994} state explicitly that in that setting the greedy algorithm
and KTNS are \emph{no longer valid}. What survives is the linear-programming
formulation. The matrix remains an interval matrix and therefore totally unimodular,
so the loading problem is still solvable in polynomial time, by solving the linear
relaxation. This distinction matters here because the Benders subproblem of
Section~\ref{ssec:bbc} is that linear program, not the KTNS rule, and so remains
exact in a regime where the rule does not apply.

A later study \citep{Crama2007} draws the boundary of tractability sharply. Once
tools may occupy several magazine slots, the loading problem alone becomes strongly
NP-hard, by reduction from $3$-partition, and yet for a fixed capacity it remains
polynomial through a shortest-path model. The uniform SSP studied here, whose loading
problem is the paging problem of Remark~\ref{rem:belady} \citep{Belady1966}, sits at
the tractable edge of a family that becomes intractable as soon as the tools stop
being uniform.

\subsection{The compact formulations, and where they stop}\label{ssec:sota-compact}

The exact side of the field is a two-decade effort to give the ordering problem a
linear relaxation strong enough to drive branch and bound. Its trajectory is the
single most important piece of context for this report.

\citet{Laporte2004} replaced the original Tang--Denardo model, and it is their paper
that established why a replacement was needed: they exhibited a feasible solution of
the Tang--Denardo linear relaxation with objective value $0$, valid on every
instance. Their own formulation, written here as LSS, introduces a dummy job as a
start-and-end depot and uses arc variables, magazine-state variables and switch
variables on a travelling-salesman skeleton. They strengthen it with three families
of valid inequalities, and report that on the Tang--Denardo example the relaxation
rises from $0$ to $6.0$ against an integer optimum of $7.0$.

They then solve the model two ways. A linear-programming branch-and-cut, seeded with
the GENIUS heuristic of \citet{Hertz1998}, and a branch-and-bound that carries no
linear relaxation at all but two combinatorial bounds on partial sequences, in the
enumerative spirit later developed by \citet{Yanasse2009}. Their computational tables
report the branch-and-cut on instances with eight jobs, and the branch-and-bound
reaching twenty-five. That a search with no relaxation beats a search with one is the
first concrete evidence that on this problem the bound, not the search, is the
limiting factor.

\citet{Ghiani2010} recast the SSP as a least-cost Hamiltonian cycle problem with a
nonlinear arc cost. They establish a symmetry property: a job sequence and its
reversal have the same total switch cost. They are careful to note that this does not
make the problem symmetric, since individual transition costs do differ; it is the
total that is preserved. The property halves the search space, and they combine it
with comb inequalities and lower bounds recomputed inside the search.

\citet{Catanzaro2015} produced the tightest compact models to date, a family of five
formulations in arc-based variables. Two of their results are used here. Formulations
3 and 4 have the same linear relaxation, and Formulation 4 is the smaller of the two;
their computational section says so explicitly and excludes Formulation 3 on that
ground. Formulation 5 has a relaxation at least as strong as Formulation 4, and their
experiments confirm it is strictly stronger on most of their datasets. Their
conclusions recommend Formulation 5 as the best integer-programming formulation for
the SSP. The present study uses Formulation 4 as its compact baseline, because that
is what was implemented; Formulation 5 was not implemented, and
Section~\ref{ssec:limits} records this as a limitation rather than a preference.

\citet{daSilva2024} give a multicommodity-flow model in which each tool routes one
unit of flow through a position network. It requires no lazy constraints, so a
standard solver can be used directly, and it includes a symmetry-breaking constraint
that eliminates half of the symmetric solutions. Its relaxation is what makes it
important. They prove that the linear relaxation of their model equals $M-C$ in
their notation: the number of tools minus the magazine capacity. Their argument
assumes, as they state, that all $M$ tools are required by some job. Under that
assumption $M=\lvert\U\rvert$, and the bound is exactly the coverage bound $q$ of
Proposition~\ref{prop:lb2}. They report solving $17.01\%$ more instances than the
models of \citet{Tang1988}, \citet{Laporte2004} and \citet{Catanzaro2015}, on the
instance sets of \citet{Yanasse2009}.

That is where the compact models reviewed above stop: the strongest relaxation among
them coincides with a counting argument that fits in one sentence. \citet{daSilva2024}
also record, for the earlier models, that their relaxations fall below $M-C$ in the
great majority of cases; this is a statement about most instances, not all, and
Table~\ref{tab:methods} in Section~\ref{sec:exact} is hedged accordingly.

\begin{table}[t]\centering
\scriptsize
\setlength{\tabcolsep}{4pt}
\begin{tabular}{@{}lll p{5.9cm}@{}}
\toprule
Year & Work & Approach & Contribution used here \\
\midrule
1988 & \citet{Tang1988}      & ILP; the KTNS rule           & KTNS optimal for a fixed order; the pairwise bound $w_{ij}$ \\
1994 & \citet{Crama1994}     & complexity; heuristics       & NP-hard for every fixed $b\ge2$; the seven standard assumptions; KTNS invalid under non-uniform tool costs, the LP is not \\
2004 & \citet{Laporte2004}   & the LSS model; B\&C and B\&B & the Tang--Denardo relaxation is identically $0$; B\&B to $25$ jobs beats B\&C at $8$ \\
2010 & \citet{Ghiani2010}    & Hamiltonian cycle            & reversal symmetry preserves the total cost \\
2015 & \citet{Catanzaro2015} & compact F1--F5               & F3 and F4 share a relaxation, F4 is smaller; F5 is strongest and is their recommendation \\
2021 & \citet{daSilva2024}   & multicommodity flow          & relaxation $=M-C$ when every tool is used; $+17.01\%$ instances solved \\
2024 & \citet{Akhundov2024}  & symmetry; arc flow           & symmetry-broken position and flow models \\
2025 & \citet{Legrand2025}   & dynamic programming, $A^{*}$ & combinatorial bounds computed inside the search \\
\bottomrule
\end{tabular}
\caption{The exact-method lineage for the SSP. Two decades of compact formulations
raised the linear relaxation only as far as the coverage bound $q=\lvert\U\rvert-b$
\citep{daSilva2024}. Moving past it has required leaving the compact
linear-programming setting.}
\label{tab:lineage}
\end{table}

\subsection{The configuration view}\label{ssec:sota-config}

A parallel line abandons job positions for tool configurations.

\citet{Crama1994Column} formulate the Job Grouping Problem as a set-covering model
solved by column generation. Its relaxation is strong, and its pricing subproblem is
a Set-Union Knapsack Problem \citep{Goldschmidt1994}, which is NP-hard even under
restrictive conditions. That difficulty returns as the bottleneck of the pricing
problems in Section~\ref{ssec:pos}. \citet{Jans2013} give a compact
asymmetric-representatives alternative, which is what the grouping solver used in
this work implements.

The configuration view that unifies grouping and sequencing is developed by
\citet{Colares2026Exact}, work in progress in the advisors' own group. It casts the
SSP as a generalised travelling-salesman problem whose clusters overlap, and solves
it by a non-compact integer program with exponentially many transition variables,
together with generalised subtour-elimination constraints separated as lazy
constraints. The overlap is what separates the problem from the textbook generalised
travelling-salesman problem with disjoint clusters surveyed by \citet{Pop2024}, whose
Noon--Bean reduction \citep{Noon1993} and branch-and-cut methods
\citep{Fischetti1997} therefore do not transfer directly. Where a pricing problem
generates constraints that depend on the columns already generated, the
column-and-row framework of \citet{Muter2013} is the relevant methodology.

Two recent theses sit closest to this work, and both come from the same research
group. \citet{Otiai2024} builds an exact branch-and-price for the Job Grouping
Problem using Ryan--Foster branching \citep{Ryan1981}. \citet{Moreira2026} studies
the SSP through the magazine-configuration formulation and measures its relaxation
against the compact models. His measurements are directly comparable with those of
Section~\ref{sec:comp}, and one of them is quoted there: he reports that the trivial
bound $M-b$ obtained by the multicommodity model was the best available bound on all
thirty instances of three of his subgroups, while on the tighter subgroup the
configuration model gave the best bound on four instances.

\subsection{Heuristics and metaheuristics}\label{ssec:sota-heur}

Because the exact methods scale poorly, most of the SSP literature is heuristic.
Early work pairs a constructive travelling-salesman-style rule with local search over
the job order \citep{Crama1994}. The GENI and GENIUS heuristics of \citet{Hertz1998}
sharpen those constructions and outperform the earlier ones; GENIUS also supplies the
upper bound inside the branch-and-cut of \citet{Laporte2004}. Enumerative schemes with
dominance pruning \citep{Yanasse2009} and the bounding arguments of
\citet{Ghiani2010} complete the pre-metaheuristic period.

The current state of the art is the hybrid genetic search of \citet{Mecler2021},
which couples an order-based genetic algorithm with local search and seeds its
population from grouping solutions. The grouping view is therefore built into the
best heuristic as well as into the exact methods.

\subsection{The decomposition in industrial practice}\label{ssec:sota-practice}

The group-then-sequence decomposition is not only a theoretical construct. In colour
printing the SSP appears in exactly the form studied here, with ink cartridges as
tools and the cartridge rack as the magazine. \citet{Burger2015} solve the industrial
problem by this decomposition, modelling the grouping as a unicost set-covering
problem and the sequencing as a travelling-salesman problem, and enumerating optimal
groupings by Stirling-number arguments.

Their assessment of the decomposition is the closest thing in the literature to the
question Section~\ref{sec:gap} attacks, and it is conditional. They concur with the
earlier finding that the decomposition does well, but qualify it: the agreement holds
\emph{especially in cases where the largest job size is only slightly smaller than
the magazine size}, and, as the difference between the magazine size and the largest
job size grows, the quality of the solutions the heuristic finds decreases. That
qualification matches the structural analysis of Section~\ref{sec:gap} closely, and
Section~\ref{ssec:disc-agree} returns to the agreement.

Their paper also poses, as an open question in its closing paragraph, the trade-off
between minimising the number of tools changed and the number of times tools are
changed, for varying values of the setup time incurred when the machine is stopped.
Section~\ref{ssec:collapse} answers that question for the uniform model.

What the literature does not contain is a worst-case analysis of the decomposition:
bounds on its gap or ratio, extremal instances, or conditions for exactness. The
survey of \citet{Calmels2018} classifies sixty-one SSP papers and identifies gaps in
the field; among them is the observation that few authors propose decomposition
methods for the SSP at all.

The exception, and the closest prior work to Section~\ref{sec:gap}, is
\citet{CramaKlundert1999}. They study the worst-case behaviour of approximation
algorithms for tool management problems, and their Theorem~8 relates the job grouping
and tool switching problems in both directions: from any grouping into batches one
obtains a schedule costing at most $b$ times the grouping value, because at most $b$
switches are needed between consecutive batches. That is precisely the construction
analysed here, and it makes the decomposition a $b$-approximation. What their analysis
does not do is sharpen the ratio below $b$, characterise the instances on which the
construction is exact, or distinguish the configuration-walk reading of the method from
the value obtained after its final loading step --- and
Section~\ref{ssec:frames} shows that the last distinction changes the answer to the
first two.

\subsection{Recent exact developments}\label{ssec:sota-recent}

The exact frontier has moved again recently, by routes that reinforce the reading
above. \citet{Akhundov2024} exploit the reversal symmetry of \citet{Ghiani2010}
systematically, deriving symmetry-broken position-based and arc-flow formulations
from a job-group representation. \citet{Legrand2025}, with Catanzaro among the
authors, leave integer programming altogether for a dynamic-programming and $A^{*}$
search equipped with a new family of lower bounds, reporting the strongest exact
results to date.

Their bounds are combinatorial and are computed inside the search rather than as a
linear relaxation. That is the pattern. Every recent advance has come from leaving
the compact linear-programming setting, whether for a search with sharper
combinatorial bounds, for a decomposition whose subproblem is settled exactly, or for
a relaxation engineered to exceed the counting bound. The two methods of
Section~\ref{sec:exact} take the second and third of those routes.

\subsection{Benders decomposition}\label{ssec:sota-benders}

The solver of Section~\ref{ssec:bbc} inherits a developed methodology.
\citet{Rahmaniani2017} survey the acceleration of Benders decomposition. The
particular setting in which the subproblem is settled by a combinatorial oracle
rather than by linear-programming duality is the logic-based Benders of
\citet{Hooker2003} and the combinatorial Benders cuts of \citet{Codato2006}. Two
classical strengthenings are used later: the Pareto-optimal cut of
\citet{MagnantiWong1981}, which selects among alternative dual optima the cut that
cuts deepest at a fixed interior point, and the practical form of
\citet{Papadakos2008}, which removes the extra constrained solve the original method
required.

The immediate source, however, is not a paper but an internship. \citet{Kashou2025}
applies branch-and-Benders-cut to a robust production-line balancing problem,
integrating Benders into the branch-and-bound tree through lazy-constraint callbacks,
replacing a general solver in the dual subproblem with four tailored algorithms, and
strengthening the method with a heuristic that pre-injects cuts. He reports that even
so the decomposition does not systematically beat the non-decomposed formulation. It
was that report which prompted Prof.~Colares to propose the same architecture for the
SSP, on the reasoning that the SSP occupies the regime the assembly-line problem
lacked: its subproblem is not a general program needing tailored algorithms, but a
polynomially solvable oracle, so the cut is exact and cheap to produce at every
callback.

\subsection{What is open, and what this report attacks}\label{ssec:sota-open}

Three observations set the agenda.

First, the two-phase heuristic has a strong empirical record, one conditional
industrial assessment \citep{Burger2015}, and no worst-case theory.
Section~\ref{sec:gap} builds one.

Second, on the exact side the strongest of the compact relaxations reviewed here equals
the coverage bound $q$ \citep{daSilva2024}, and the configuration relaxation does not improve on it on
most instances \citep{Moreira2026}. Progress therefore requires either a method that
does not rely on a linear relaxation at all, which is the Benders route of
Section~\ref{ssec:bbc}, or a formulation whose relaxation can exceed the bound, which
is the position--transition route of Section~\ref{ssec:pos}.

Third, the configuration view that works well for the Job Grouping Problem
\citep{Otiai2024} has no tractable exact counterpart for the SSP, because the natural
model carries exponentially many subtour constraints. That is the opening the
position-indexed formulations address.
